\chapter{Introduction and Background}
\label{ch:intro}

% BUDGET: 7 pages (was 6; took one from ch3 for the positioning section).

\section{Motivation}
\label{sec:motivation}

Academic research contains a great deal of work that is necessary but not itself scholarly.
Finding and triaging the relevant literature, translating an idea into an experiment,
implementing that experiment correctly, and assembling the result into a structured document
consume time that is not spent on the judgement that constitutes the research. This is
familiar enough to be unremarkable. What makes it worth addressing is that the burden is not
distributed evenly.

The mechanical parts of modern research are disproportionately \emph{programming} parts. A
researcher who cannot write code is not merely slower at running an experiment; they are
frequently unable to run it at all without recruiting somebody who can. The requirements
analysis conducted for this project, reported in Chapter~\ref{ch:requirements}, illustrates
the point directly: asked what they found most frustrating about designing and running
experiments, one respondent answered simply, \emph{``my own programming abilities''}, and went
on to say that experiment code was the one stage of the research process they would willingly
delegate to a machine. The same respondent refused to delegate literature review or writing,
regarding those as constitutive of scholarship rather than overhead. That distinction ---
between the parts of research that are the work and the parts that merely stand in its way ---
is the distinction this project attempts to serve.

The prospect of using large language models to close that gap is not new, and several systems
discussed in Section~\ref{sec:background} already attempt it. What has not been adequately
addressed is the condition under which such a system could responsibly be deployed inside a
university. Two obstacles stand out. The first is trust: a system that produces a plausible
literature review, a plausible experiment and a plausible set of results gives a
non-programming user no means of telling a correct result from a fabricated one, and the
scale of that fabrication is now measured rather than suspected. Auditing 111 million
references across 2.5 million papers, Zhao et al.~\cite{zhao2026hallucinations} find a sharp
rise in non-existent citations following widespread adoption of language models, with a
conservative estimate of 146,932 hallucinated citations in 2025 alone, and report that
existing safeguards --- preprint moderation and journal peer review --- capture only a
fraction of them. The second is governance: an
unpublished hypothesis, an unpublished result and an early draft are among the most sensitive
artefacts a research institution holds, and sending them to a third-party commercial API is
a disclosure that many institutions cannot sanction.

This project therefore sets out to build a research assistant intended for institutional
deployment: one that runs on the university's own hardware under its own control, and whose
outputs are constructed so that a person who cannot read the code it wrote can nevertheless
establish whether to believe them.

\section{Problem statement}
\label{sec:problem}

\noindent\textbf{Input:} A research question, expressed in plain text.

\noindent\textbf{Task:} Retrieve and synthesise the relevant literature; derive testable
hypotheses from it; plan a corresponding experiment; generate and execute the code for that
experiment; and produce a structured research paper reporting what was found --- such that
every factual claim in the resulting document is mechanically traceable to evidence the system
actually obtained, and any claim that is not so traceable is withheld rather than asserted.

The final clause is the difficult one, and it is what separates this project from a pipeline
that simply chains model calls together. A system may be permitted to be wrong; it may not be
permitted to be confidently, unverifiably wrong to a user who has no means of checking it.

\section{Background}
\label{sec:background}

\subsection{Autonomous research agents}

The immediate predecessor to this work is Agent Laboratory~\cite{schmidgall2025agentlab},
which accepts a human research idea and progresses through literature review, experimentation
and report writing, producing a code repository and a written report. It decomposes the
pipeline into agents named after academic roles --- PhD Student, Postdoc, ML Engineer,
Professor --- and supplies specialised sub-systems (\emph{mle-solver}, \emph{paper-solver})
for code generation and report assembly. Two of its reported findings shaped this project more
than its architecture did. First, human involvement through a ``co-pilot'' mode materially
improves output quality relative to fully autonomous operation. Second, and more troubling,
its automated reviewer substantially overestimates the quality of the papers it assesses
relative to human reviewers. A system that grades its own work generously is not merely
inaccurate; it is inaccurate in the direction that most misleads a user who is relying on it.

The AI Scientist~\cite{lu2024aiscientist} pursues fully autonomous end-to-end discovery,
including simulated peer review, and ResearchAgent~\cite{baek2024researchagent} iteratively
refines research ideas using reviewing agents. Both depend on proprietary models. In each of
these systems the reviewing component is itself a language model asked for a judgement, which
is precisely the arrangement that produces the overestimation Agent Laboratory reports.

The wider field is broader than the general-purpose pipelines this work sits among. Agentic
systems have been applied within single disciplines, at a scale and with a degree of physical
integration well beyond anything attempted here --- MARS~\cite{shi2026mars} coordinates
nineteen agents and sixteen domain tools with laboratory robotics for materials discovery ---
and to narrower slices of the research process, such as academic writing
support~\cite{jain2026researchhelper}. Yu et al.~\cite{yu2024ai4r} survey the area as an
emerging research paradigm in its own right.

\subsection{The two systems closest to this work}

Two recent systems bear directly on what is claimed here, and the relationship to each is
worth stating precisely rather than leaving to inference.

\textbf{Autonomous Research Loops}~\cite{koyun2026loops} is the closest comparable system: an
LLM-agent framework covering hypothesis generation, literature search, coding, execution,
analysis, manuscript preparation and peer-style review, built on lightweight experiment
templates, producing a paper for roughly fifteen dollars. Its architecture and this one were
arrived at independently and resemble each other closely. What differs is the response to
failure. Koyun reports three major failure modes --- hallucination of experimental details,
broken implementations, and misinterpretation of metrics --- and concludes that such pipelines
hold promise \emph{``provided that effective verification mechanisms and safety constraints are
incorporated from the outset''}. Those three failures are precisely what
Chapter~\ref{ch:design} builds mechanisms against, and the verification that paper identifies
as necessary is what this project contributes. Its reviewing component is itself a language
model performing multi-pass self-critique, which is the arrangement Agent
Laboratory~\cite{schmidgall2025agentlab} found to overestimate quality --- so the concern
raised in Section~\ref{sec:background} is now visible in two independent systems.

\textbf{AutoMETA}~\cite{lee2026autometa} is the closest precedent for this project's central
design principle. Its agents extract outcomes from individual papers and then hand the
finalised dataset \emph{``to a non-agent statistical module''} for pooling, and its authors
conclude that procedural enforcement yields reliable statistical reasoning \emph{``without
relying solely on model capability''}. An ablation supports the claim: removing inter-agent
critique and protocol enforcement degrades median effect-size error from 6.4\% to 12.4\%,
25.1\% and 28.1\%. This is independent corroboration, from a different domain, that moving
work out of the model improves reliability --- and it establishes the principle as prior art
rather than a claim original to this dissertation.

The distinction is nonetheless a real one. AutoMETA moves a \emph{computation} out of the
model: a pooled effect size is arithmetic, and arithmetic belongs in code. This project moves
\emph{verification} out: whether a citation refers to a paper that was actually retrieved,
whether a verdict follows from code that actually ran, whether reported data was real. The
first replaces model judgement where a correct answer is computable; the second constrains
model output where it is not. AutoMETA is also confined to a single well-defined task, whereas
the pipeline here spans retrieval to written paper, which is what makes the number of places
requiring such a constraint the substance of the contribution.

\subsection{Verification and originality}

Two independent bodies of evidence establish that outputs from systems of this class cannot
be taken at face value.

The first concerns citation. Zhao et al.~\cite{zhao2026hallucinations} audit 111 million
references across arXiv, bioRxiv, SSRN and PubMed Central, and find hallucinated citations
concentrated in exactly the conditions this project operates under: fields with rapid uptake
of language models, manuscripts bearing the linguistic signature of machine assistance, and
small or early-career author teams. Two of their findings bear directly on the design that
follows. They observe that automated detection pipelines built on language models
\emph{``might reintroduce model-specific biases into the analysis''} --- an argument against
checking a model's output with another model, which is the arrangement most systems in this
area adopt. And they find that hallucinated references disproportionately credit already
prominent scholars, so the failure is not merely factual but distributive; that consequence is
taken up in Chapter~\ref{ch:ethics}.

The second concerns novelty. Gupta and Pruthi~\cite{gupta2025glitters} provide the strongest
available evidence on originality. In an expert-led study of
fifty research documents produced by state-of-the-art research agents, 24\% contained clear
plagiarism --- defined as one-to-one methodological mapping or significant uncredited
borrowing --- and every plagiarised document had passed the systems' own built-in novelty
checks. Automated detection tools, including Turnitin, identified none of the verified
instances. The authors further observe that documents which expert reviewers had rated
\emph{more novel} than human-written work were subsequently found to be plagiarised once those
reviewers were specifically directed to look for overlap.

Two conclusions follow, and both are load-bearing here. Evaluating such a system on the
apparent quality of its output measures the wrong thing, because the failure mode is
invisible to a reader not actively hunting for it. And a checking mechanism implemented as a
model judging its own output does not work, because the same model that produced the
overlap is being asked to notice it.

\subsection{Graph-structured orchestration}

The system described in this dissertation is built on
LangGraph~\cite{langgraph}, which models a computation as a directed graph of nodes over a
shared, checkpointed state. The relevant property is not the programming convenience but the
consequence for inspection and recovery: every step, including deterministic non-model steps,
is an addressable node whose inputs and outputs are recorded, so a run can be resumed from
where it stopped and examined stage by stage after the fact. Chapter~\ref{ch:design} explains
where that mattered and where it did not.

\subsection{Open-weight models and institutional deployment}

The systems above depend on commercial APIs. This project uses open-weight models served on
the university's own high-performance computing facility, principally Qwen3-Coder
30B-A3B-Instruct~\cite{qwen3coder} served through vLLM~\cite{kwon2023vllm}. The motivation is
not cost. It is that an institutional research tool must not export its users' unpublished
work, and self-hosting is the only arrangement in which that guarantee is structural rather
than contractual. A consequence, developed throughout Chapter~\ref{ch:design}, is that the
system must remain correct while being driven by a model considerably weaker than the
commercial frontier --- which turns out to sharpen the design rather than compromise it.

\section{Aims and requirements}
\label{sec:aims}

\subsection{Aims}

\begin{description}[style=unboxed,leftmargin=0pt,itemsep=0.3em]
  \item[\textbf{A1}] Enable a researcher to proceed from a research question to a
    literature-grounded, experimentally-supported draft paper \emph{without writing code}.
  \item[\textbf{A2}] Ensure that no claim the system presents is unverifiable: citations must
    resolve to papers actually retrieved, hypothesis verdicts must derive from code that
    actually ran, and results computed on synthesised inputs must withhold a conclusion rather
    than assert one.
  \item[\textbf{A3}] Keep every stage independently usable and steerable, so that a researcher
    may use only the parts they need and override the system's choices rather than accept them.
  \item[\textbf{A4}] Run entirely on institutional infrastructure using open-weight models, so
    that unpublished research does not leave the university.
  \item[\textbf{A5}] Make generated experiments recover from their own failures without human
    intervention, since the intended user cannot debug them.
  \item[\textbf{A6}] Survive the conditions of a shared HPC facility --- pre-emption, queue
    limits and storage quotas --- so that an unattended run is resumable rather than lost.
\end{description}

\subsection{Functional requirements}

\begin{description}[style=unboxed,leftmargin=0pt,itemsep=0.2em,font=\normalfont\bfseries]
  \item[FR1] Accept a plain-text research question as the only mandatory input. \hfill(A1)
  \item[FR2] Retrieve, deduplicate and download relevant literature from open-access
    scholarly APIs. \hfill(A1)
  \item[FR3] Identify adjacent research fields and surface concrete cross-disciplinary
    connections to the question. \hfill(A1)
  \item[FR4] Synthesise the retrieved literature into testable, grounded hypotheses, ranked
    against one another. \hfill(A1)
  \item[FR5] Convert a hypothesis into an implementation-ready experiment plan, with an
    explicit feasibility judgement. \hfill(A1)
  \item[FR6] Generate runnable experiment code and, where the environment permits, execute it
    and capture structured results. \hfill(A1, A5)
  \item[FR7] Diagnose failures in generated code and regenerate it, bounded by an explicit
    attempt budget. \hfill(A5)
  \item[FR8] Draft a structured research paper from the resulting evidence, rendered as a PDF.
    \hfill(A1)
  \item[FR9] Review that paper against the same upstream evidence the drafting stage used, and
    revise it in response. \hfill(A2)
  \item[FR10] Permit a user to run a selected range of stages, supply their own papers, and
    choose which hypothesis is carried forward. \hfill(A3)
  \item[FR11] Present live progress, intermediate outputs and retrieved sources through a
    graphical interface requiring no terminal. \hfill(A1, A3)
\end{description}

\subsection{Non-functional requirements}

\begin{description}[style=unboxed,leftmargin=0pt,itemsep=0.2em,font=\normalfont\bfseries]
  \item[NFR1] Every decision that can be settled mechanically must be settled in code rather
    than by model judgement. \hfill(A2)
  \item[NFR2] No research content may be transmitted to a third-party model provider.
    \hfill(A4)
  \item[NFR3] Each stage must validate its input against the preceding stage's published
    contract before acting on it. \hfill(A3)
  \item[NFR4] A run interrupted by pre-emption or crash must be resumable from the stage at
    which it stopped. \hfill(A6)
  \item[NFR5] All non-model logic must be unit-testable without a language model present.
    \hfill(A2)
\end{description}

\subsection{Divergence from the original project proposal}
\label{sec:divergence}

The project proposal for this work specified a reinforcement learning component as its
principal novelty: agent behaviour was to be improved over iterations using reward signals
derived from code execution outcomes and paper scores, targeting hallucination, poor
self-evaluation and instruction-following failures. That component was not delivered in the
form proposed. During implementation it became apparent that all three of its stated targets
were better addressed by moving decisions out of the model altogether than by optimising the
model's behaviour within it, and the effort was redirected accordingly.
Chapter~\ref{ch:po} states that argument in full, quotes the proposal's own justification
against it, and discharges the requirement with a preference dataset derived from the
system's own failure history together with a direct preference optimisation
experiment~\cite{rafailov2023dpo}.

Three smaller divergences are recorded here for completeness and discussed where relevant.
A similarity-checking stage, promised in response to Gupta and Pruthi's
findings~\cite{gupta2025glitters}, was not implemented; it emerged from the requirements
analysis as the highest-priority requirement by participant rating, and its absence is the
project's most significant shortfall (Section~\ref{sec:future-overlap}). Reports are rendered
directly to PDF rather than through LaTeX, for reasons given in Chapter~\ref{ch:design}.
Conversely, both features the proposal listed as optional --- a graphical interface and a
co-pilot mode --- were delivered, and substantially exceeded.

\section{Contributions}
\label{sec:contributions}

\begin{description}[style=unboxed,leftmargin=0pt,itemsep=0.25em,font=\normalfont\bfseries]
  \item[C1] A seven-agent research pipeline in which every agent is itself a graph, not merely
    the orchestrator, so that each individual model call and each deterministic check is an
    independently addressable and inspectable step.
  \item[C2] A systematic application of the principle that verifiable decisions belong in code
    rather than in the model, across every stage of a full research pipeline. The principle
    itself is not original --- AutoMETA~\cite{lee2026autometa} demonstrates it for statistical
    pooling, and this work was arrived at independently rather than derived from it. What is
    contributed here is its extension from computation to \emph{verification}, and the
    catalogue of the places an end-to-end pipeline turns out to need it, each paired with the
    concrete failure that motivated it.
  \item[C3] A reliability architecture for language-model-generated experiment code:
    classification of failures by kind before response, reduction of the cost of a failed
    attempt before its frequency, separate budgets for malformed and merely incorrect
    responses, and regeneration confined to the code region actually implicated.
  \item[C4] A data provenance mechanism that withholds a hypothesis verdict when an
    experiment's inputs were synthesised, on the grounds that reporting a refutation computed
    on invented data is a worse outcome than reporting nothing.
  \item[C5] An empirical failure analysis of the system in operation, including a documented
    case in which a feature intended to improve grounding actively degraded the trustworthiness
    of results.
  \item[C6] A preference dataset constructed from the system's own fix loop, in which each
    rejected generation, the check that rejected it, and its accepted replacement form a
    preference triple by construction.
\end{description}

\section{Terminology}
\label{sec:terminology}

The reader is assumed to be able to program and to be familiar with language models in
general terms. The following usages are specific to this project.

\begin{description}[style=unboxed,leftmargin=0pt,itemsep=0.2em]
  \item[\emph{Agent}] A self-contained pipeline stage with a published input and output
    contract, implemented internally as a graph. The term carries no implication of a
    persona; Chapter~\ref{ch:requirements} explains why role-based naming was abandoned.
  \item[\emph{Verdict}] The classification of a hypothesis as supported, refuted or
    inconclusive. Computed in code from execution results, never asserted by the model.
  \item[\emph{Surrogate data}] Data synthesised by generated code because no real source for a
    declared input could be resolved. Any surrogate input causes the verdict to be withheld.
  \item[\emph{Fix loop}] The cycle in which failing generated code is diagnosed, regenerated
    against the concrete error, and re-executed, bounded by an explicit attempt budget.
  \item[\emph{Smoke run}] A preliminary execution of generated code with every cost parameter
    reduced to its minimum, used to surface defects in seconds rather than after a full-length
    run.
  \item[\emph{Bridge insight}] A concrete statement that a method from an adjacent field could
    inform the research question, citing the cross-field papers it was drawn from.
\end{description}

\section{Structure of this dissertation}

Chapter~\ref{ch:requirements} reports the requirements analysis conducted with academic staff
before implementation, and the design decisions traceable to it. Chapter~\ref{ch:design}
presents the design and implementation of the system together, covering the architecture, the
verification mechanisms that give effect to A2, the reliability work behind A5, the deployment
arrangements behind A4 and A6, and the approaches that were tried and abandoned.
Chapter~\ref{ch:po} addresses the reinforcement learning requirement discussed in
Section~\ref{sec:divergence}. Chapter~\ref{ch:evaluation} reports the testing regime, a
quantitative benchmark of the code-generation stage, an analysis of the system's characteristic
failures, the human evaluation, and an assessment against A1--A6.
Chapter~\ref{ch:ethics} sets out the project's ethical position, Chapter~\ref{ch:conclusion}
concludes and identifies further work, and Chapter~\ref{ch:bcs} addresses the BCS project
criteria and reflects critically on the project as a whole.
